\begin{textsample}{POS Dim 1 – se – Score 32.00 – se/se\_ef\_3.txt}  \label{ex:f1_pos_115}
EDUCAÇÃO FÍSICA As reformas curriculares têm se tornado expressões de processos de mudanças na história da educação do Brasil.
Essa é a tendência \textbf{que} vem se repetindo, sistematicamente, de acordo com cada momento do desenvolvimento econômico, cultural e político.
Este currículo se destina aos professores do Ensino Fundamental.
Foi desenvolvido a partir da construção coletiva entre professores das Diretorias de Educação da Secretaria Estadual de Educação ( SEED/SE ), das Secretarias Municipais de Educação de Sergipe e do Colégio de Aplicação da Universidade Federal de Sergipe ( Codap/UFS ), em regime de colaboração.
Deve nortear a elaboração dos Planos de Curso e dos Planos de Aula para a \textbf{disciplina} Educação Física, como também dialogar com os Projetos Político-Pedagógicos de cada escola.
Ressalte -se \textbf{que} caberá ao professor construir seu próprio percurso, de acordo com os \textbf{sujeitos} atendidos e a realidade da escola a \textbf{que} pertence, para atingir \textbf{um} objetivo comum, garantido na Base Nacional Comum Curricular, \textbf{que} é proporcionar aos alunos o acesso aos conhecimentos da cultura corporal.
No componente curricular Educação Física existe várias concepções educacionais e pedagógicas \textbf{que} defendem preceitos ideológicos diferentes.
Antes de se abordar quaisquer aspectos acerca dos elementos \textbf{constitutivos} dessa \textbf{disciplina} no currículo sergipano, há necessidade de se esclarecer \textbf{que}, para ser implementada \textbf{enquanto} prática pedagógica de forma qualificada, precisa ser legitimada numa perspectiva vinculada às necessidades e interesses daqueles \textbf{que} serão envolvidos no processo de ensino e de aprendizagem.
Assim, deve -se tomar como referência princípios \textbf{norteadores} \textbf{que} se fundamentam a partir de orientações \textbf{teórico-metodológicas} e tem como base \textbf{uma} formação \textbf{que} visa à emancipação dos \textbf{sujeitos}.
É nesse contexto \textbf{que} vislumbramos a Educação Física \textbf{enquanto} componente curricular, sendo \textbf{uma} fonte de conhecimentos, habilidades e atitudes \textbf{que} possibilitem a concepção de \textbf{um} \textbf{sujeito} \textbf{que} se reconhece como cidadão detentor de papel crítico e construtivo dentro do seu micro e macro espaço político, social, tecnológico, cultural e econômico.
Neste sentido, espera -se \textbf{que} este componente curricular possa estabelecer espaços de reflexões e diálogo junto aos educandos acerca da cultura corporal, compreendida como elementos \textbf{historicamente} construídos pelo \textbf{homem}.
Pensando dessa forma, a Educação Física propiciará discussões e reflexões sobre a corporeidade[5 ], mostrando os sentidos e significados \textbf{que} cada corpo assumiu \textbf{historicamente} na \textbf{singularidade} e na diversidade sociocultural.
Dentro do cenário brasileiro altamente contraditório desigual e excludente, destaca -se a região nordeste e suas mazelas, conjuntura na qual a educação física é atraída a adotar \textbf{uma} postura efetivamente crítica, tendo por objetivo a construção de \textbf{uma} prática pedagógica balizada em elementos \textbf{que} proponham a transformação dessa realidade a partir de vivências \textbf{que}, no âmbito da cultura corporal, \textbf{sistematizem} de forma progressiva a \textbf{assimilação} ativa dos conhecimentos, habilidades e atitudes, e medeiem junto com a \textbf{ludicidade} à construção de \textbf{um} comportamento crítico e emancipado, contribuindo para o discernimento sobre a relevância da organização individual e coletiva, do respeito, da igualdade e da \textbf{equidade} no seu relacionamento com os outros, dentre outras questões.
Acerca das questões culturais, é importante considerar a resistência, a diferença e a luta por visibilidade e reconhecimento \textbf{que} habitam o interior de \textbf{um} mesmo grupo cultural.
Sendo a escola \textbf{um} espaço \textbf{constituinte} da teia social e \textbf{um} dos primeiros ambientes onde o contato ( e o conflito ) entre os diferentes se manifesta, fazendo emergir as lutas e movimentos de resistência cada vez mais \textbf{explícitos}, é inviável tomá -los por irrelevantes ou dissimular sua existência.
Tal fato enuncia a problemática dos confrontos identitários na escola multicultural. [ 5 ] : Propriedade \textbf{que} nos garante a compreensão do corpo.
Para Santini, 1987, é \textbf{uma} expressão usada pela filosofia para definir a \textbf{percepção} do ser humano como indivisível, sem a \textbf{divisão} corpo e \textbf{mente}.
Na escola democrática destes tempos, \textbf{uma} educação multiculturalmente orientada \textbf{implica} a assunção de \textbf{uma} postura clara em favor da luta contra a opressão, o preconceito e a discriminação aos quais foram submetidos alguns grupos \textbf{historicamente} desprovidos de poder, sem \textbf{que} se \textbf{perca} de vista a perene composição de novos grupos culturais ( NEIRA ; NUNES, 2009, p. 210 ).
Ao considerar a hegemonia dos valores de \textbf{uma} cultura branca e masculina, o currículo multicultural crítico visa a abrir espaço para o estudo de práticas corporais pertencentes tanto aos diversos grupos étnicos como às mulheres, analisar criticamente o predomínio das manifestações da cultura corporal branca e as masculinas, dentre outras manifestações \textbf{que} guardem relações assimétricas de poder, como por exemplo, a de classe social.
Nesse sentido, o currículo deve abrir espaços para \textbf{que} os rappers e skatistas estudem melhor o rap e o skate e também as demais práticas corporais, enfim, para \textbf{que} estudem o afoxé, o funk etc., sempre acompanhado das histórias de luta desses movimentos, pelo seu reconhecimento e \textbf{dignidade} ( NEIRA, 2007, p. 158 ).
Isso significa \textbf{que}, na Educação Física, deve -se estudar as diferentes manifestações da cultura corporal de forma contextualizada, acompanhada pela participação de representantes dos grupos \textbf{que} as recriam, desenvolvem e praticam.
Nada \textbf{que} se assemelhe à perspectiva de “ visitas ” descontextualizadas a diferentes manifestações da cultura corporal.
Todos os grupos inseridos nas instituições escolares têm a necessidade de se unir na luta comum pelo fortalecimento da democracia.
Nesse sentido, propõe -se buscar \textbf{uma} prática \textbf{que} entenda a natureza específica da diferença, mas \textbf{que} também aprecie a adesão comum aos princípios de igualdade e justiça.
A intenção é reconhecer e problematizar as \textbf{categorias} \textbf{que} constroem as representações das identidades para \textbf{que} os estudantes possam compreender os significados das diferenças \textbf{que} separam os interesses dos indivíduos de grupos diversos e se engajem no respeito e valorização da diversidade e de suas identidades.
Nesse sentido é \textbf{que} propomos, no currículo sergipano, inserir objetos de aprendizagem e especialmente habilidades \textbf{que} tocam nas questões socioculturais, na diversidade, na identificação da diferença, na \textbf{problematização} das relações de poder e tencionam \textbf{uma} formação mais ampla e crítica para o estudante sergipano no percurso com a Educação Física.
Atrelados a essas questões, os conteúdos/conhecimentos têm função importantíssima e podem ser entendidos como a produção humana construída \textbf{historicamente} e \textbf{que}, vinculados ao currículo escolar, devem retratar a interatividade do \textbf{homem} com a natureza e com o meio social.
Ressalte -se \textbf{que} o currículo formal é envolvido por relações de poder \textbf{demarcando} espaços e discursos.
Sendo assim, deve ser compreendido para além da listagem de matérias, do rol de \textbf{disciplinas}, mas como algo dinâmico, \textbf{um} movimento da escola \textbf{que} constrói \textbf{uma} base material capaz de realizar o projeto de escolarização do \textbf{homem}, construída por três aspectos : o trato com o conhecimento, a organização escolar e a normatização escolar.
A Lei de Diretrizes e Bases da Educação Nacional ( LDB - Lei n.o 9.394/1996 ) estabeleceu, em seu artigo 26, os rumos \textbf{que} a Educação Física deveria seguir, promulgada na tentativa de transformar o ensino brasileiro.
Com o lançamento dos Parâmetros Curriculares Nacionais em 1997, surgiu a possibilidade de elaborar \textbf{um} programa curricular integrado à proposta pedagógica da escola visando a ampliar o conhecimento do aluno, justificando a Educação Física \textbf{enquanto} componente curricular da Educação Básica.
Já a Base Nacional Comum Curricular ( BNCC ) garante o \textbf{que} todo estudante deve aprender nas escolas brasileiras ao longo da educação básica.
Nela a Educação Física está localizada na área de linguagens e é o componente curricular \textbf{que} tematiza as práticas corporais ( Jogos e Brincadeiras, Lutas, Esportes, Ginásticas, Práticas corporais de Aventura e Danças ) em suas diversas formas de codificação e significação social, entendidas como manifestações das possibilidades \textbf{expressivas} dos \textbf{sujeitos} e patrimônio cultural da \textbf{humanidade} ( BNCC, 2017 ).
Ressalta -se \textbf{que} apesar de não ter sido apresentada como \textbf{umas} das práticas corporais organizadas do componente curricular, é necessário \textbf{que} os estudantes sergipanos tenham a oportunidade de vivenciar práticas no meio líquido, pois além de contribuir para o desenvolvimento de habilidades motoras e valências físicas, constitui a identidade cultural do nosso povo, \textbf{visto} \textbf{que} estamos localizado no nordeste do Brasil, região de clima quente, com várias cidades \textbf{que} contam com a presença de praias e rios \textbf{que} montam o cenário de \textbf{uma} gama de práticas corporais de lazer e esportivas.
Assim, a Base Nacional Comum Curricular reitera : Essa afirmação não se vincula apenas à ideia de vivenciar e/ou aprender, por exemplo, os esportes aquáticos ( em especial, a natação em seus quatro estilos competitivos ), mas também à proposta de experimentar “ atividades aquáticas ”.
São, portanto, práticas centradas na ambientação dos estudantes ao meio líquido \textbf{que} permitem aprender, entre outros movimentos básicos, o controle da respiração, a flutuação em equilíbrio, a imersão e os deslocamentos na água ( BNCC, 2017, p. 217 ).
O estudo desse conhecimento visa aprender a expressão corporal como linguagem dando caráter crítico a esse conteúdo ( SOARES \textbf{et} al., 2012 ). “ Nas aulas, as práticas corporais devem ser abordadas como fenômeno cultural dinâmico, diversificado, pluridimensional, singular e contraditório ” ( BNCC, 2017, p. 211 ).
Deve -se dar destaque \textbf{que} todas as práticas corporais vivenciadas na escola devem ser reconstruídas com base em sua função social e suas possibilidades materiais.
Isso significa \textbf{dizer} \textbf{que} tais práticas podem ser transformadas no interior da escola.
Por exemplo, as práticas corporais de aventura devem ser adaptadas às condições da escola, ocorrendo de maneira simulada, tomando -se como referência o cenário de cada contexto escolar.
Mas também é importante o engajamento de professores para explorar os espaços existentes no entorno da escola, no bairro, na cidade e no Estado \textbf{que} permitem a realização dessas práticas de modo real, na vida cultural dos \textbf{sujeitos}.
É relevante ressaltar \textbf{que} os trabalhos com as práticas corporais de aventura possam e devam oportunizar ao aluno, reflexões sobre a preservação ambiental e sustentabilidade.
A organização das unidades temáticas se \textbf{baseia} na compreensão de \textbf{que} o caráter lúdico está presente em todas as práticas corporais, ainda \textbf{que} esta não seja a finalidade da Educação Física na escola.
Ao brincar, dançar, jogar, praticar esportes, ginásticas ou atividades de aventura, para além da \textbf{ludicidade}, os estudantes se apropriam das lógicas intrínsecas ( regras, códigos, rituais, sistemáticas de funcionamento, organização, táticas etc. ) a essas manifestações, assim como trocam entre si e com a sociedade as representações e os significados \textbf{que} lhes são atribuídos.
Por essa razão, a delimitação das habilidades privilegia oito dimensões de conhecimento : Experimentação ; Uso e Apropriação ; Fruição ; Reflexão sobre a ação ; Construção de valores ; Análise ; Compreensão e Protagonismo comunitário.

% matched lemmas: assimilação, basear, categoria, constituinte, constitutivo, demarcar, dignidade, disciplina, divisão, dizer, enquanto, equidade, et, explícito, expressivo, historicamente, homem, humanidade, implicar, ludicidade, mente, norteador, percepção, perder, problematização, que, singularidade, sistematizar, sujeito, teórico-metodológico, um, ver
\end{textsample}
